\documentclass[DIN, pagenumber=false, parskip=half,
               fromalign=right, fromphone=false, fromfax=false, frombank=true,
               yourmail=false,
               yourref=false,
               fromrule=false]{scrlttr2}


%\usepackage{ngerman}
\usepackage[utf8]{inputenc}
\usepackage[T1]{fontenc}
\usepackage{textcomp}
\usepackage{longtable}
\usepackage{eurosym}

\LoadLetterOption{template}

\newcommand{\stdlohn}[0]{11 \euro}

\renewcommand*{\raggedsignature}{\raggedright}

\setkomavar{fromname}{GNUnet e. V.}
\setkomavar{fromaddress}{Boltzmannstraße 3 \\ 85748 Garching bei München}
%\setkomavar{frombank}{IBAN: DE67830654080004822650 \\ BIC: GENODEF1SLR}

\setkomavar{subject}[]{Recommendation for nlnet for the EC's NGI initiative}
%\setkomavar{yourmail}[Ihre Angebots-Nr.]{2011012631}
%\setkomavar{yourref}[Ihre Kundennummer]{263}
\setkomavar{date}[Datum]{\today}


%===================================
% footer
\firstfoot{
	\parbox[t]{\textwidth}{\footnotesize

	\begin{tabular}[t]{l@{}}
		 \multicolumn{1}{@{}l@{}}{Association:}\\
		 GNUnet e. V. \\
		 Boltzmannstraße 3, 85748 Garching bei München \\ \\
		 E-Mail: ev@gnunet.org \\
	\end{tabular}
	\hfill{}
	\begin{tabular}[t]{l@{}}
	        \multicolumn{1}{@{}l@{}}{\usekomavar*{frombank}:}\\
	        \usekomavar{frombank}
	\end{tabular}
 }}

\begin{document}


%===================================
% receiver
\begin{letter}{
        To whom it may concern \\
}

\opening{\ }
\vspace{-1.5cm}

%===================================

GNUnet e.V. is an association dedicated to supporting research,
development and education in the area of secure decentralized
networking.  We are excited to learn that NLnet will apply for funding
from the EC's NGI initiative, and are happy to lend them our strongest
support.

Various community members of GNUnet e.V. have worked with NLnet over
the last decade, and we have found their funding process to be
reasonably uncomplicated and effective.  In terms of knowledgeable and
result-oriented mechanisms for grantmaking, NLnet is comparable with
the (less specialized) Emmy Noether Programme of the DFG, except that
it's decision making speed is even faster.  We also have experience
with H2020, FP7, KTI and NSF funding proceedures, which are all more
complicated and generally result in lower quality expert reviews (of
both applications and deliverables).  We can also attest that NLnet is
well known in the larger community of people building modern,
privacy-preserving network applications and services, so we predict
that they will have no problem attracting a substantial number of
quality applications.

We thus think that it is a really good idea for NLnet to coordinate
funding in this area: many promising projects are unable to directly
meet the requirements to receive EU funding, in particular given the
complex application and reporting procedures.  Furthermore, NLnet's
plan to create a shared infrastructure would add some of the
professionalism that is currently lacking in this space.  As GNUnet
e.V., we see particularly strong value in infrastructure to address
usability, accessibility, continuous integration and reproducible
builds, as these are not currently adequately covered for us by the
GNU project.

Thus, as one of the larger and older efforts in the space of
reinventing the Internet, GNUnet e.V. would really like to see NLnet
to be responsible for managing the NGI and selecting the right
(sub)projects.  Having them as an intermediary will be hugely
beneficial as we know that they can adapt to the needs of the
people that might actually make a difference.


\vspace{3cm}
Prof. Dr. Christian Grothoff\\
President, GNUnet e. V. \\
{\texttt grothoff@gnunet.org}


\end{letter}
\end{document}
